% =============================================================================
% APPENDIX GI: LIT-GIGERENZER: Ecological Rationality and Fast-and-Frugal Heuristics
% =============================================================================
% Module: BCM2_LIT_GIGERENZER
% Category: LIT
% Version: 1.0 (January 2026)
% Status: Final
%
% Comprehensive integration of Gerd Gigerenzer's research contributions
% on fast-and-frugal heuristics, ecological rationality, risk communication,
% and the adaptive toolbox into the EBF framework.
%
% Language: English
% =============================================================================

% =============================================================================
% CROSS-REFERENCE STRUCTURE
% =============================================================================
%
% CHAPTER LINKAGE (Required):
% - Primary Chapter: Chapter 11: Awareness (CORE-AWARE)
% - Secondary Chapters: Chapter 9: Context, Chapter 3: Limits of Classical Utility
%
% This appendix integrates with:
%
% DEPENDENCIES (what this appendix REQUIRES):
% - Appendix L (LIT-KAHNEMAN): Contrast with heuristics-and-biases
% - Appendix V (CORE-WHEN): Context framework (ecological structure)
% - Appendix AU (CORE-AWARE): Awareness function
% - Appendix BBB (CORE-WHERE): Parameter repository
%
% DEPENDENTS (what REQUIRES this appendix):
% - Appendix BBB (CORE-WHERE): Parameters for LLMMC
% - Appendix AH (DOMAIN-CONSULTING): Practical decision tools
% - All DOMAIN appendices using heuristic decision rules
%
% =============================================================================

\begin{tcolorbox}[colback=gray!5!white, colframe=gray!75!black,
    title=Appendix Metadata]
\textbf{Appendix:} GI --- LIT-GIGERENZER \\
\textbf{Category:} Literature Integration \\
\textbf{Author Focus:} Gerd Gigerenzer (Max Planck Institute, Berlin) \\
\textbf{Papers:} 19 (Mixed Tier Evidence) \\
\textbf{Primary Contribution:} Fast-and-frugal heuristics, ecological rationality, risk literacy
\end{tcolorbox}

\begin{tcolorbox}[colback=purple!5!white, colframe=purple!75!black,
    title=Chapter Linkage]

\textbf{Primary Chapter:} Chapter 11: Awareness
\begin{itemize}[nosep]
\item Section 11.3: Heuristic Processing $\leftrightarrow$ This Appendix Section \ref{sec:gi-heuristics}
\item Section 11.8: Information Processing $\leftrightarrow$ This Appendix Section \ref{sec:gi-ecological}
\end{itemize}

\textbf{Secondary Chapters:}
\begin{itemize}[nosep]
\item Chapter 9: Context --- ecological structure as $\Psi$
\item Chapter 3: Limits of Classical Utility --- bounded rationality foundations
\end{itemize}

\textbf{Reading Guide:}
\begin{itemize}[nosep]
\item For conceptual overview: Read Chapter 11 first
\item For contrast: Compare with Appendix L (LIT-KAHNEMAN)
\item For formal details: This appendix provides complete specification
\end{itemize}

\end{tcolorbox}


\chapter{LIT-GIGERENZER: Ecological Rationality and Fast-and-Frugal Heuristics}
\label{app:lit-gigerenzer}


\begin{abstract}
This appendix integrates Gerd Gigerenzer's research program into the EBF framework.
Gigerenzer's work establishes that: (1) simple heuristics can outperform complex
optimization in uncertain environments (less-is-more); (2) rationality should be
evaluated by fit between decision strategy and environment structure (ecological
rationality); (3) risk communication improves with natural frequencies over
probabilities. His contributions provide an alternative to the ``bias'' framing,
showing when heuristics are adaptive rather than erroneous. 19 papers spanning
1987-2011 are integrated, offering a complementary perspective to Kahneman-Tversky.
\end{abstract}


\begin{tcolorbox}[colback=blue!5!white, colframe=blue!75!black,
    title=Quick Reference: LIT-GIGERENZER]

\textbf{Core Question:} When and why do simple heuristics outperform complex optimization?

\textbf{Key Formula (Take-the-Best):}
\begin{equation}
\text{Choose } a \text{ if } \exists i: c_i(a) > c_i(b) \text{ and } \forall j < i: c_j(a) = c_j(b)
\end{equation}
where $c_i$ are cues ordered by validity.

\textbf{Key Concepts:}
\begin{itemize}[nosep]
\item \textbf{Fast-and-frugal heuristics}: Simple rules using limited information
\item \textbf{Ecological rationality}: Match between heuristic and environment
\item \textbf{Less-is-more}: More information can reduce accuracy
\item \textbf{Adaptive toolbox}: Multiple heuristics for different problems
\end{itemize}

\textbf{Cross-References:} Appendix L (LIT-KAHNEMAN), V (CORE-WHEN), AU (CORE-AWARE), BBB (CORE-WHERE)
\end{tcolorbox}

% -----------------------------------------------------------------------------
% APPENDIX SCOPE BOX
% -----------------------------------------------------------------------------

\begin{tcolorbox}[
    colback=orange!5!white,
    colframe=orange!75!black,
    title={\textbf{Appendix Scope: Ziel / In-Scope / Out-of-Scope / Constraints}},
    fonttitle=\bfseries
]

\textbf{Ziel (Objective):}
\begin{itemize}[nosep]
    \item When does using less information and simpler rules lead to better decisions than optimization with full information?
\end{itemize}

\textbf{In-Scope:}
\begin{itemize}[nosep]
    \item Theoretical Framework
    \item Axioms and Formal Definitions
    \item Key Results and Findings
    \item Theme 1: Fast-and-Frugal Heuristics
    \item Theme 2: Ecological Rationality
\end{itemize}

\textbf{Out-of-Scope (delegiert an andere Appendices):}
\begin{itemize}[nosep]
    \item LIT-KAHNEMAN $\rightarrow$ Appendix L
    \item CORE-WHEN $\rightarrow$ Appendix V
    \item CORE-AWARE $\rightarrow$ Appendix AU
    \item CORE-WHERE $\rightarrow$ Appendix BBB
    \item CORE-WHERE $\rightarrow$ Appendix BBB
\end{itemize}

\textbf{Constraints (Anwendungsgrenzen):}
\begin{itemize}[nosep]
    \item [TODO: Define when this appendix does NOT apply]
    \item [TODO: Define required prerequisites]
    \item [TODO: Define known limitations]
\end{itemize}

\textbf{Lieferobjekte (Deliverables):}
\begin{itemize}[nosep]
    \item 7 Table(s)
    \item 16 Equation(s)
    \item 6 Axiom(s)
    \item Worked Example(s)
\end{itemize}

\end{tcolorbox}




\begin{tcolorbox}[colback=red!5!white, colframe=red!75!black,
    title=The Fundamental Question]
\textbf{When does using less information and simpler rules lead to better decisions than optimization with full information?}
\end{tcolorbox}


%%%%%%%%%%%%%%%%%%%%%%%%%%%%%%%%%%%%%%%%%%%%%%%%%%%%%%%%%%%%%%%%%%%%%%%%%%%%%%%
\section{Theoretical Framework}
\label{sec:gi-theory}
%%%%%%%%%%%%%%%%%%%%%%%%%%%%%%%%%%%%%%%%%%%%%%%%%%%%%%%%%%%%%%%%%%%%%%%%%%%%%%%

Gigerenzer's research program offers a fundamentally different perspective from the heuristics-and-biases tradition:

\subsection{The Rationality Debate}

While Kahneman-Tversky document deviations from normative models as ``biases,'' Gigerenzer asks: \textit{biased relative to what?}

\begin{equation}
\text{Rationality}_{\text{KT}} = \frac{\text{Behavior}}{\text{Normative Model}} \quad \text{vs.} \quad
\text{Rationality}_{\text{GI}} = \frac{\text{Behavior}}{\text{Environment Structure}}
\end{equation}

\subsection{The Ecological Turn}

Gigerenzer argues that heuristics should be evaluated by their fit to environmental structure:
\begin{equation}
\text{Performance}(h, E) = f(\text{Match}(h, E))
\end{equation}
where $h$ is a heuristic and $E$ is the environment.

\subsection{The Less-is-More Effect}

Counter-intuitively, using less information can \textit{improve} predictions:
\begin{equation}
\text{Accuracy}(h_{\text{simple}}) > \text{Accuracy}(h_{\text{complex}}) \quad \text{when } E \text{ has high uncertainty}
\end{equation}


%%%%%%%%%%%%%%%%%%%%%%%%%%%%%%%%%%%%%%%%%%%%%%%%%%%%%%%%%%%%%%%%%%%%%%%%%%%%%%%
\section{Axioms and Formal Definitions}
\label{sec:gi-axioms}
%%%%%%%%%%%%%%%%%%%%%%%%%%%%%%%%%%%%%%%%%%%%%%%%%%%%%%%%%%%%%%%%%%%%%%%%%%%%%%%

Based on Gigerenzer's empirical findings, we formalize six axioms for the EBF:

\begin{tcolorbox}[colback=gray!5!white, colframe=gray!75!black,
    title=Axiom GI-1: Less-is-More Effect]
\textbf{Statement:} In environments with high uncertainty, simple heuristics can outperform complex optimization.
\begin{equation}
\text{If } \sigma^2_E > \sigma^2_{\text{crit}}: \quad \mathbb{E}[\text{Accuracy}(h_{\text{frugal}})] > \mathbb{E}[\text{Accuracy}(h_{\text{optimal}})]
\end{equation}
\textbf{Interpretation:} Overfitting to noise is avoided by ignoring information.

\textbf{Source:} \citet{PAP-gigerenzer1999simple}, \citet{PAP-gigerenzer2009homo}
\end{tcolorbox}

\begin{tcolorbox}[colback=gray!5!white, colframe=gray!75!black,
    title=Axiom GI-2: Ecological Rationality]
\textbf{Statement:} Heuristic performance depends on the match between heuristic structure and environment structure.
\begin{equation}
\text{Performance}(h, E) = g(\text{Structure}(h), \text{Structure}(E))
\end{equation}
\textbf{Interpretation:} The same heuristic can be brilliant in one environment and terrible in another.

\textbf{Source:} \citet{gigerenzer2001bounded}, \citet{PAP-goldstein2002recognition}
\end{tcolorbox}

\begin{tcolorbox}[colback=gray!5!white, colframe=gray!75!black,
    title=Axiom GI-3: Recognition Heuristic]
\textbf{Statement:} If one of two objects is recognized and the other is not, infer that the recognized object has the higher criterion value.
\begin{equation}
\text{If } R(a) = 1 \text{ and } R(b) = 0: \quad \text{Predict } a > b
\end{equation}
\textbf{Interpretation:} Recognition is a valid cue when recognition correlates with the criterion.

\textbf{Source:} \citet{PAP-goldstein2002recognition}
\end{tcolorbox}

\begin{tcolorbox}[colback=gray!5!white, colframe=gray!75!black,
    title=Axiom GI-4: Take-the-Best]
\textbf{Statement:} Search through cues in order of validity; stop search when a cue discriminates; choose the option favored by that cue.
\begin{equation}
\text{Decision} = \arg\max_{x} c_{i^*}(x) \quad \text{where } i^* = \min\{i : c_i(a) \neq c_i(b)\}
\end{equation}
\textbf{Interpretation:} Lexicographic decision making that ignores trade-offs.

\textbf{Source:} \citet{PAP-gigerenzer1996fast}, \citet{PAP-brandstatter2006priority}
\end{tcolorbox}

\begin{tcolorbox}[colback=gray!5!white, colframe=gray!75!black,
    title=Axiom GI-5: Natural Frequency Superiority]
\textbf{Statement:} Bayesian reasoning improves when probabilities are presented as natural frequencies.
\begin{equation}
\text{Accuracy}_{\text{frequency}} > \text{Accuracy}_{\text{probability}}
\end{equation}
\textbf{Interpretation:} Minds evolved to process frequencies, not single-event probabilities.

\textbf{Source:} \citet{PAP-gigerenzer1995improve}, \citet{PAP-gigerenzer2012helping}
\end{tcolorbox}

\begin{tcolorbox}[colback=gray!5!white, colframe=gray!75!black,
    title=Axiom GI-6: Adaptive Toolbox]
\textbf{Statement:} Humans possess a repertoire of heuristics and select among them based on environmental cues.
\begin{equation}
h^*(E) = \arg\max_{h \in \mathcal{H}} \text{Performance}(h, E)
\end{equation}
\textbf{Interpretation:} No single heuristic is universally best; the toolbox is adaptive.

\textbf{Source:} \citet{gigerenzer2001bounded}, \citet{PAP-gigerenzer2008rationality}
\end{tcolorbox}


%%%%%%%%%%%%%%%%%%%%%%%%%%%%%%%%%%%%%%%%%%%%%%%%%%%%%%%%%%%%%%%%%%%%%%%%%%%%%%%
\section{Key Results and Findings}
\label{sec:gi-results}
%%%%%%%%%%%%%%%%%%%%%%%%%%%%%%%%%%%%%%%%%%%%%%%%%%%%%%%%%%%%%%%%%%%%%%%%%%%%%%%

\begin{table}[htbp]
\centering
\caption{Summary of Key Results from Gigerenzer's Research}
\label{tab:gi-key-results}
\renewcommand{\arraystretch}{1.3}
\begin{tabular}{@{}llcc@{}}
\toprule
\textbf{Result} & \textbf{Finding} & \textbf{Effect} & \textbf{Tier} \\
\midrule
GI-1 & Take-the-best matches regression & Equal accuracy & 2 \\
GI-2 & Recognition heuristic is valid & $r > 0.6$ with criterion & 2 \\
GI-3 & Frequency formats improve reasoning & +40\% accuracy & 1 \\
GI-4 & Less-is-more in prediction & Simpler $>$ complex & 2 \\
GI-5 & Ecological match predicts success & Strong correlation & 2 \\
GI-6 & Heuristics beat optimization OOS & Cross-validation & 2 \\
GI-7 & Gut feelings can be accurate & Domain-dependent & 5 \\
GI-8 & Biases disappear with frequencies & 65\% $\to$ 90\% & 1 \\
\bottomrule
\end{tabular}
\end{table}


%%%%%%%%%%%%%%%%%%%%%%%%%%%%%%%%%%%%%%%%%%%%%%%%%%%%%%%%%%%%%%%%%%%%%%%%%%%%%%%
\section{Theme 1: Fast-and-Frugal Heuristics}
\label{sec:gi-heuristics}
%%%%%%%%%%%%%%%%%%%%%%%%%%%%%%%%%%%%%%%%%%%%%%%%%%%%%%%%%%%%%%%%%%%%%%%%%%%%%%%

\subsection{The Core Insight}

Gigerenzer's foundational work \citep{PAP-gigerenzer1999simple} establishes that simple heuristics are not merely ``good enough'' approximations---they can \textit{outperform} optimization:

\begin{itemize}[nosep]
\item \textbf{Fast}: Require little computation
\item \textbf{Frugal}: Use little information
\item \textbf{Accurate}: Often match or beat complex models
\end{itemize}

\subsection{The Recognition Heuristic}

The recognition heuristic \citep{PAP-goldstein2002recognition} is the simplest:
\begin{quote}
If you recognize one option but not the other, choose the recognized option.
\end{quote}

\paragraph{When it works.} Recognition is ecologically valid when:
\begin{itemize}[nosep]
\item Recognition correlates with criterion (e.g., city size correlates with media mention)
\item One option is recognized, the other is not (discrimination possible)
\end{itemize}

\paragraph{Empirical demonstration.} German students outperformed American students in predicting which US city is larger---because Germans recognized fewer cities, they could use the recognition heuristic more effectively.

\subsection{Take-the-Best}

Take-the-best \citep{PAP-gigerenzer1996fast} is a lexicographic heuristic:

\begin{enumerate}[nosep]
\item Order cues by validity (predictive power)
\item Search cues in order
\item Stop when a cue discriminates between options
\item Choose the option favored by that cue
\end{enumerate}

\paragraph{Key property.} Take-the-best ignores trade-offs---it never compensates for a negative on one cue with positives on others. This is ``irrational'' by compensatory utility standards but can be optimal in sparse environments.

\subsection{Implications for EBF}

Fast-and-frugal heuristics suggest that CORE-AWARE should model:
\begin{equation}
A_{\text{eff}}(\Psi) = \begin{cases}
A_{\text{heuristic}}(\Psi) & \text{if } \Psi \in \mathcal{E}_{\text{favorable}} \\
A_{\text{deliberative}}(\Psi) & \text{if } \Psi \in \mathcal{E}_{\text{unfavorable}}
\end{cases}
\end{equation}
where $\mathcal{E}_{\text{favorable}}$ are environments where heuristics match.


%%%%%%%%%%%%%%%%%%%%%%%%%%%%%%%%%%%%%%%%%%%%%%%%%%%%%%%%%%%%%%%%%%%%%%%%%%%%%%%
\section{Theme 2: Ecological Rationality}
\label{sec:gi-ecological}
%%%%%%%%%%%%%%%%%%%%%%%%%%%%%%%%%%%%%%%%%%%%%%%%%%%%%%%%%%%%%%%%%%%%%%%%%%%%%%%

\subsection{The Ecological Framework}

Ecological rationality \citep{gigerenzer2001bounded} shifts the normative question from ``Does behavior match the optimal model?'' to ``Does behavior match the environment?''

\begin{equation}
\text{Ecological Rationality} = \text{Fit}(\text{Mental Strategy}, \text{Environmental Structure})
\end{equation}

\subsection{Environmental Structure Types}

Different environments favor different heuristics:

\begin{table}[htbp]
\centering
\caption{Environment-Heuristic Matching}
\label{tab:gi-environments}
\renewcommand{\arraystretch}{1.3}
\begin{tabular}{@{}lll@{}}
\toprule
\textbf{Environment} & \textbf{Favored Heuristic} & \textbf{Why} \\
\midrule
High uncertainty & Simple heuristics & Avoid overfitting \\
One dominant cue & Take-the-best & Exploit cue structure \\
Recognition valid & Recognition heuristic & Use ecological correlation \\
Time pressure & Fast heuristics & Computation costs \\
Low predictability & Equal weighting & No cue is reliable \\
\bottomrule
\end{tabular}
\end{table}

\subsection{Connection to CORE-WHEN ($\Psi$)}

Gigerenzer's environmental structure maps to EBF's context framework:
\begin{equation}
\Psi_{\text{ecological}} = (\text{Uncertainty}, \text{Cue Structure}, \text{Time Pressure}, \text{Feedback Quality})
\end{equation}


%%%%%%%%%%%%%%%%%%%%%%%%%%%%%%%%%%%%%%%%%%%%%%%%%%%%%%%%%%%%%%%%%%%%%%%%%%%%%%%
\section{Theme 3: Risk Communication}
\label{sec:gi-risk}
%%%%%%%%%%%%%%%%%%%%%%%%%%%%%%%%%%%%%%%%%%%%%%%%%%%%%%%%%%%%%%%%%%%%%%%%%%%%%%%

\subsection{The Natural Frequency Revolution}

Gigerenzer's risk communication research \citep{gigerenzer1995improve, gigerenzer2007helping} demonstrates that Bayesian reasoning failures disappear when probabilities are replaced with natural frequencies:

\paragraph{Probability format (difficult):}
\begin{quote}
``The probability of breast cancer is 1\%. If a woman has breast cancer, the probability of a positive mammogram is 90\%. If she does not have breast cancer, the probability of a positive mammogram is 9\%. What is the probability that a woman with a positive mammogram actually has breast cancer?''
\end{quote}
\textit{Result: Only 15\% of physicians answer correctly (correct answer: $\approx$10\%).}

\paragraph{Natural frequency format (easy):}
\begin{quote}
``10 out of 1,000 women have breast cancer. 9 of these 10 women will have a positive mammogram. Of the 990 women without cancer, 89 will have a positive mammogram. How many of the women with positive mammograms actually have cancer?''
\end{quote}
\textit{Result: 65\% of physicians answer correctly.}

\subsection{Implications}

Natural frequencies improve reasoning because:
\begin{itemize}[nosep]
\item They preserve base rate information
\item They make the reference class explicit
\item They match how statistical information was historically encoded (frequencies, not single-event probabilities)
\end{itemize}


%%%%%%%%%%%%%%%%%%%%%%%%%%%%%%%%%%%%%%%%%%%%%%%%%%%%%%%%%%%%%%%%%%%%%%%%%%%%%%%
\section{Theme 4: The Kahneman-Gigerenzer Debate}
\label{sec:gi-debate}
%%%%%%%%%%%%%%%%%%%%%%%%%%%%%%%%%%%%%%%%%%%%%%%%%%%%%%%%%%%%%%%%%%%%%%%%%%%%%%%

\subsection{The Tension}

Gigerenzer and Kahneman represent two different research programs on bounded rationality:

\begin{table}[htbp]
\centering
\caption{Kahneman vs. Gigerenzer: Two Programs}
\label{tab:gi-kt-debate}
\renewcommand{\arraystretch}{1.3}
\begin{tabular}{@{}lll@{}}
\toprule
\textbf{Aspect} & \textbf{Kahneman-Tversky} & \textbf{Gigerenzer} \\
\midrule
Focus & Biases and errors & Adaptive heuristics \\
Normative standard & Logic, probability theory & Ecological fit \\
Heuristics are... & Sources of systematic error & Tools that can excel \\
Research style & Lab experiments, puzzles & Computational modeling \\
Policy implication & Debias, nudge & Improve environments \\
\bottomrule
\end{tabular}
\end{table}

\subsection{EBF Integration}

EBF integrates both perspectives:
\begin{itemize}[nosep]
\item \textbf{Kahneman-Tversky}: Documents when heuristics fail (important for intervention design)
\item \textbf{Gigerenzer}: Documents when heuristics succeed (important for not over-correcting)
\end{itemize}

The synthesis for CORE-AWARE:
\begin{equation}
A(\Psi) = \begin{cases}
\text{Heuristic-appropriate} & \text{if } \text{Match}(\Psi, h) > \theta \\
\text{Bias-prone} & \text{if } \text{Match}(\Psi, h) < \theta
\end{cases}
\end{equation}


%%%%%%%%%%%%%%%%%%%%%%%%%%%%%%%%%%%%%%%%%%%%%%%%%%%%%%%%%%%%%%%%%%%%%%%%%%%%%%%
\section{Integration with Other Appendices}
\label{sec:gi-integration}
%%%%%%%%%%%%%%%%%%%%%%%%%%%%%%%%%%%%%%%%%%%%%%%%%%%%%%%%%%%%%%%%%%%%%%%%%%%%%%%

\subsection{Connection to Appendix L (LIT-KAHNEMAN)}

\begin{tcolorbox}[colback=yellow!5!white, colframe=yellow!75!black,
    title=Integration: LIT-GIGERENZER $\leftrightarrow$ LIT-KAHNEMAN]
\textbf{What LIT-GIGERENZER provides:}
\begin{itemize}[nosep]
\item Environmental conditions where heuristics succeed
\item Less-is-more as counterpoint to bias research
\item Ecological validity perspective
\end{itemize}

\textbf{What LIT-KAHNEMAN provides:}
\begin{itemize}[nosep]
\item Environmental conditions where heuristics fail
\item Systematic bias documentation
\item Intervention targets (nudges)
\end{itemize}

\textbf{EBF Synthesis:} Match heuristic to environment; intervene when mismatch.
\end{tcolorbox}

\subsection{Connection to Appendix V (CORE-WHEN)}

Gigerenzer's ecological structure enriches $\Psi$:
\begin{equation}
\Psi = (\Psi_E, \Psi_S, \Psi_T, ..., \Psi_{\text{ecological}})
\end{equation}
where $\Psi_{\text{ecological}}$ captures cue structure, uncertainty level, and feedback quality.

\subsection{Connection to Appendix AU (CORE-AWARE)}

The awareness function must account for heuristic-environment match:
\begin{equation}
A_{\text{Gigerenzer}} = A_{\text{base}} \cdot \text{Match}(h, \Psi)
\end{equation}


%%%%%%%%%%%%%%%%%%%%%%%%%%%%%%%%%%%%%%%%%%%%%%%%%%%%%%%%%%%%%%%%%%%%%%%%%%%%%%%
\section{Worked Example: Designing Risk Communication}
\label{sec:gi-example}
%%%%%%%%%%%%%%%%%%%%%%%%%%%%%%%%%%%%%%%%%%%%%%%%%%%%%%%%%%%%%%%%%%%%%%%%%%%%%%%

\subsection{Setup}
A health organization wants to communicate cancer screening risks. Using Gigerenzer's framework to design effective communication.

\paragraph{The Problem.} Patients misunderstand screening test results, leading to over-treatment and anxiety.

\paragraph{Application with Natural Frequencies.}

\textbf{Step 1: Identify the statistics to communicate.}
\begin{itemize}[nosep]
\item Sensitivity: 90\% (true positive rate)
\item Specificity: 91\% (true negative rate)
\item Base rate: 1\% (prevalence)
\end{itemize}

\textbf{Step 2: Convert to natural frequencies.}
Instead of: ``If you test positive, there's an X\% chance you have cancer''

Use: ``Imagine 1,000 people like you are tested:
\begin{itemize}[nosep]
\item 10 have cancer; 9 of these test positive
\item 990 don't have cancer; 89 of these test positive
\item Of 98 positive tests, only 9 have cancer''
\end{itemize}

\textbf{Step 3: Make the comparison explicit.}
``About 1 in 10 positive tests indicates actual cancer.''

\paragraph{Result.}
Natural frequency format:
\begin{itemize}[nosep]
\item Physicians: Accuracy increases from 15\% to 65\%
\item Patients: Better calibrated expectations
\item Policy: Reduced unnecessary biopsies
\end{itemize}

\paragraph{Discussion.}
This example shows how EBF can use Gigerenzer's framework to design interventions that work \textit{with} cognitive architecture rather than against it.


%%%%%%%%%%%%%%%%%%%%%%%%%%%%%%%%%%%%%%%%%%%%%%%%%%%%%%%%%%%%%%%%%%%%%%%%%%%%%%%
\section{Implications for Practice}
\label{sec:gi-practice}
%%%%%%%%%%%%%%%%%%%%%%%%%%%%%%%%%%%%%%%%%%%%%%%%%%%%%%%%%%%%%%%%%%%%%%%%%%%%%%%

\begin{tcolorbox}[colback=green!5!white, colframe=green!75!black,
    title=Practical Implications]
\begin{enumerate}
\item \textbf{Match heuristics to environments:} Don't assume optimization is always better; simple rules can win
\item \textbf{Use natural frequencies:} Present statistical information as frequencies, not probabilities
\item \textbf{Exploit ecological correlations:} Recognition and other cues can be valid; don't dismiss them as biases
\item \textbf{Design for the adaptive toolbox:} People have multiple heuristics; help them select the right one
\item \textbf{Less can be more:} In high-uncertainty environments, simpler models may predict better
\end{enumerate}
\end{tcolorbox}


%%%%%%%%%%%%%%%%%%%%%%%%%%%%%%%%%%%%%%%%%%%%%%%%%%%%%%%%%%%%%%%%%%%%%%%%%%%%%%%
\section{Parameters for LLMMC Estimation}
\label{sec:gi-llmmc}
%%%%%%%%%%%%%%%%%%%%%%%%%%%%%%%%%%%%%%%%%%%%%%%%%%%%%%%%%%%%%%%%%%%%%%%%%%%%%%%

From Gigerenzer's research, the following parameters are available for LLMMC:

\begin{table}[htbp]
\centering
\caption{LLMMC Parameters from LIT-GIGERENZER}
\label{tab:gi-llmmc}
\renewcommand{\arraystretch}{1.3}
\begin{tabular}{@{}llccc@{}}
\toprule
\textbf{Parameter} & \textbf{Description} & \textbf{Point Estimate} & \textbf{95\% CI} & \textbf{Tier} \\
\midrule
$\sigma^2_{\text{crit}}$ & Uncertainty threshold for less-is-more & 0.3 & [0.2, 0.4] & 2 \\
$v_{\text{recognition}}$ & Recognition heuristic validity & 0.65 & [0.55, 0.75] & 2 \\
$\Delta_{\text{frequency}}$ & Frequency format improvement & +0.40 & [+0.30, +0.50] & 1 \\
$n_{\text{cues}}$ & Typical cues used in take-the-best & 2-3 & [1, 4] & 2 \\
\bottomrule
\end{tabular}
\end{table}


%%%%%%%%%%%%%%%%%%%%%%%%%%%%%%%%%%%%%%%%%%%%%%%%%%%%%%%%%%%%%%%%%%%%%%%%%%%%%%%
\section{Summary}
\label{sec:gi-summary}
%%%%%%%%%%%%%%%%%%%%%%%%%%%%%%%%%%%%%%%%%%%%%%%%%%%%%%%%%%%%%%%%%%%%%%%%%%%%%%%

\begin{tcolorbox}[colback=green!10!white, colframe=green!75!black,
    title=\textbf{Appendix GI: Summary}]

\textbf{Core Question:} When do simple heuristics outperform complex optimization?

\textbf{Main Contributions:}
\begin{itemize}[nosep]
    \item Less-is-more: Simple heuristics can beat optimization in uncertain environments
    \item Ecological rationality: Evaluate strategies by environment fit, not abstract norms
    \item Fast-and-frugal heuristics: Recognition, take-the-best, 1/N
    \item Natural frequencies: +40\% improvement in Bayesian reasoning
    \item Adaptive toolbox: Multiple heuristics selected by context
    \item Kahneman-Gigerenzer synthesis: Both bias and adaptation perspectives valid
\end{itemize}

\textbf{Key Outputs:}
\begin{itemize}[nosep]
    \item Environment-heuristic matching function
    \item Less-is-more threshold ($\sigma^2_{\text{crit}} \approx 0.3$)
    \item Frequency format improvement ($\Delta \approx +40\%$)
\end{itemize}

\textbf{Integration:} This appendix connects to CORE-AWARE (heuristic processing),
CORE-WHEN (ecological structure as $\Psi$), and LIT-KAHNEMAN (complementary perspective).

\end{tcolorbox}


%%%%%%%%%%%%%%%%%%%%%%%%%%%%%%%%%%%%%%%%%%%%%%%%%%%%%%%%%%%%%%%%%%%%%%%%%%%%%%%
\section{Glossary of Symbols}
\label{sec:gi-glossary}
%%%%%%%%%%%%%%%%%%%%%%%%%%%%%%%%%%%%%%%%%%%%%%%%%%%%%%%%%%%%%%%%%%%%%%%%%%%%%%%

\begin{tcolorbox}[colback=blue!5!white, colframe=blue!75!black]
\textbf{Central Reference:} For complete symbol definitions, see
\textbf{Appendix G} (\textit{Glossary and Notation}).

This section lists only symbols introduced or specialized in this appendix.
\end{tcolorbox}

\vspace{0.5em}

\begin{table}[htbp]
\centering
\caption{Symbols Introduced in Appendix GI}
\label{tab:gi-symbols}
\renewcommand{\arraystretch}{1.3}
\begin{tabular}{@{}clp{6cm}c@{}}
\toprule
\textbf{Symbol} & \textbf{Name} & \textbf{Definition} & \textbf{Also in G?} \\
\midrule
$h$ & Heuristic & Decision rule/strategy & NEW \\
$E$ & Environment & Decision environment structure & NEW \\
$\mathcal{H}$ & Adaptive toolbox & Set of available heuristics & NEW \\
$c_i$ & Cue & Information cue with validity ranking & NEW \\
$R(x)$ & Recognition & Binary recognition indicator & NEW \\
$v_{\text{cue}}$ & Cue validity & Predictive power of cue & NEW \\
$\sigma^2_{\text{crit}}$ & Uncertainty threshold & When less-is-more applies & NEW \\
\bottomrule
\end{tabular}
\end{table}


%%%%%%%%%%%%%%%%%%%%%%%%%%%%%%%%%%%%%%%%%%%%%%%%%%%%%%%%%%%%%%%%%%%%%%%%%%%%%%%
\section{Critical Foundations and Objections}
\label{sec:gi-foundations}
%%%%%%%%%%%%%%%%%%%%%%%%%%%%%%%%%%%%%%%%%%%%%%%%%%%%%%%%%%%%%%%%%%%%%%%%%%%%%%%

\subsection{Objection 1: Less-is-More is Cherry-Picked}

\textbf{Objection:} The less-is-more effect only holds in carefully selected datasets. In most real-world settings, more information helps.

\textbf{Response:} Gigerenzer's research shows:
\begin{itemize}[nosep]
\item Less-is-more is systematic in high-uncertainty, low-sample environments
\item The effect is tied to overfitting prevention, which is well-understood statistically
\item Real-world success stories exist (e.g., hospital triage, financial forecasting)
\end{itemize}

\subsection{Objection 2: People Don't Actually Use These Heuristics}

\textbf{Objection:} The heuristics are mathematical models; there's no evidence people actually compute take-the-best.

\textbf{Response:}
\begin{itemize}[nosep]
\item Process tracing studies show cue-by-cue search patterns consistent with take-the-best
\item Response times match predictions of sequential search
\item The heuristics are ``as-if'' models with predictive validity, like other economic models
\end{itemize}

\subsection{Objection 3: This Is Just Anti-Kahneman Contrarianism}

\textbf{Objection:} Gigerenzer's program is more about criticizing Kahneman than building positive theory.

\textbf{Response:}
\begin{itemize}[nosep]
\item The adaptive toolbox is a substantive positive program
\item Ecological rationality provides a new normative framework
\item The debate has been productive for both research programs
\item EBF integrates both perspectives as complementary
\end{itemize}


%%%%%%%%%%%%%%%%%%%%%%%%%%%%%%%%%%%%%%%%%%%%%%%%%%%%%%%%%%%%%%%%%%%%%%%%%%%%%%%
\section{Open Issues and Research Agenda}
\label{sec:gi-open}
%%%%%%%%%%%%%%%%%%%%%%%%%%%%%%%%%%%%%%%%%%%%%%%%%%%%%%%%%%%%%%%%%%%%%%%%%%%%%%%

\begin{table}[htbp]
\centering
\caption{Research Roadmap for Gigerenzer-Related Questions}
\label{tab:gi-roadmap}
\renewcommand{\arraystretch}{1.3}
\begin{tabular}{@{}clcl@{}}
\toprule
\textbf{\#} & \textbf{Issue} & \textbf{Priority} & \textbf{Status} \\
\midrule
1 & Heuristic selection mechanisms & HIGH & Active research \\
2 & Boundary conditions for less-is-more & HIGH & Partial evidence \\
3 & Individual differences in toolbox & MEDIUM & Underexplored \\
4 & Training to improve heuristic use & MEDIUM & Mixed results \\
5 & Neural basis of heuristic processing & LOW & Early stage \\
\bottomrule
\end{tabular}
\end{table}


%%%%%%%%%%%%%%%%%%%%%%%%%%%%%%%%%%%%%%%%%%%%%%%%%%%%%%%%%%%%%%%%%%%%%%%%%%%%%%%
\section{Complete Paper Inventory}
\label{sec:gi-papers}
%%%%%%%%%%%%%%%%%%%%%%%%%%%%%%%%%%%%%%%%%%%%%%%%%%%%%%%%%%%%%%%%%%%%%%%%%%%%%%%

\begin{table}[htbp]
\centering
\caption{All 19 Gigerenzer Papers in EBF Database}
\label{tab:gi-papers}
\renewcommand{\arraystretch}{1.2}
\scriptsize
\begin{tabular}{@{}llll@{}}
\toprule
\textbf{ID} & \textbf{Year} & \textbf{Theme} & \textbf{Tier} \\
\midrule
gigerenzer1987cognition & 1987 & Risk Communication & 5 \\
gigerenzer1988base & 1988 & Risk Communication & 2 \\
gigerenzer1991illusions & 1991 & Heuristics & 5 \\
gigerenzer1995improve & 1995 & Natural Frequencies & 1 \\
gigerenzer1996fast & 1996 & Fast-and-Frugal & 2 \\
gigerenzer1996narrow & 1996 & KT Debate & 5 \\
gigerenzer1999simple & 1999 & Heuristics & 2 \\
todd2000precis & 2000 & Heuristics & 5 \\
gigerenzer2001bounded & 2001 & Ecological Rationality & 5 \\
goldstein2002recognition & 2002 & Recognition Heuristic & 2 \\
gigerenzer2003risks & 2003 & Risk Communication & 2 \\
gigerenzer2004mindless & 2004 & Statistics Critique & 5 \\
PAP-gigerenzer2005rain & 2005 & Risk Communication & 2 \\
brandstatter2006priority & 2006 & Priority Heuristic & 2 \\
gigerenzer2007gut & 2007 & Intuition & 5 \\
gigerenzer2007helping & 2007 & Medical Decisions & 2 \\
gigerenzer2008rationality & 2008 & Bounded Rationality & 5 \\
gigerenzer2009homo & 2009 & Homo Heuristicus & 5 \\
gigerenzer2011heuristic & 2011 & Heuristic Decision Making & 5 \\
\bottomrule
\end{tabular}
\end{table}


%%%%%%%%%%%%%%%%%%%%%%%%%%%%%%%%%%%%%%%%%%%%%%%%%%%%%%%%%%%%%%%%%%%%%%%%%%%%%%%
\section{References}
\label{sec:gi-references}
%%%%%%%%%%%%%%%%%%%%%%%%%%%%%%%%%%%%%%%%%%%%%%%%%%%%%%%%%%%%%%%%%%%%%%%%%%%%%%%

\begin{tcolorbox}[colback=blue!5!white, colframe=blue!75!black]
\textbf{Master Bibliography:} For complete reference details, see
\texttt{bcm\_master.bib} in the repository.

This section lists key references for this appendix.
\end{tcolorbox}

\vspace{0.5em}

\subsection{Key References}

\begin{itemize}[nosep]
    \item \textbf{Fast-and-Frugal Heuristics:} \citet{PAP-gigerenzer1996fast}, \citet{PAP-gigerenzer1999simple}
    \item \textbf{Recognition Heuristic:} \citet{PAP-goldstein2002recognition}
    \item \textbf{Ecological Rationality:} \citet{gigerenzer2001bounded}
    \item \textbf{Natural Frequencies:} \citet{PAP-gigerenzer1995improve}, \citet{PAP-gigerenzer2012helping}
    \item \textbf{KT Debate:} \citet{PAP-gigerenzer1996narrow}
\end{itemize}

% Link to master bibliography
\nocite{bcm_master}


%%%%%%%%%%%%%%%%%%%%%%%%%%%%%%%%%%%%%%%%%%%%%%%%%%%%%%%%%%%%%%%%%%%%%%%%%%%%%%%
% CROSS-REFERENCE FOOTER (Required)
%%%%%%%%%%%%%%%%%%%%%%%%%%%%%%%%%%%%%%%%%%%%%%%%%%%%%%%%%%%%%%%%%%%%%%%%%%%%%%%

\vspace{1em}
\hrule
\vspace{0.5em}

\noindent\textit{Cross-references:}
Appendix G (Central Glossary),
Appendix L (LIT-KAHNEMAN),
Appendix V (CORE-WHEN),
Appendix AU (CORE-AWARE),
Appendix BBB (CORE-WHERE).

\noindent\textit{Master Bibliography:} \texttt{bcm\_master.bib}


% =============================================================================
% END OF APPENDIX GI
% =============================================================================
